% Teoria do prezentacji "Gdzie mitygacja błędów przestaje się opłacać"
% Kompilacja: latexmk -pdf prezentacja_teoria.tex
\documentclass[11pt,a4paper]{article}
\usepackage[T1]{fontenc}
\usepackage[utf8]{inputenc}
\usepackage[polish]{babel}
\usepackage{lmodern}
\usepackage{microtype}
\usepackage[margin=2.4cm]{geometry}
\usepackage{amsmath,amssymb,mathtools}
\usepackage{booktabs}
\usepackage{enumitem}
\usepackage{xcolor}
\usepackage{tikz}
\usetikzlibrary{arrows.meta,positioning,calc}
\usepackage{titlesec}
\usepackage[hidelinks]{hyperref}

\definecolor{axblue}{HTML}{4AA8F5}
\definecolor{axbluedk}{HTML}{1565C0}
\definecolor{axgold}{HTML}{E3B53C}
\definecolor{axgolddk}{HTML}{8A6A0B}
\definecolor{axink}{HTML}{111318}

\setlist{itemsep=2pt,topsep=4pt}
\setlength{\parindent}{0pt}
\setlength{\parskip}{6pt}
\linespread{1.08}

\titleformat{\section}{\Large\bfseries}{\thesection}{10pt}{}
\titleformat{\subsection}{\large\bfseries}{\thesubsection}{8pt}{}
\titlespacing*{\subsection}{0pt}{14pt}{4pt}

\newcommand{\PH}{\colorbox{yellow!40}{\texttt{[\_\_]}}}
\newcommand{\slid}[1]{\hfill{\normalfont\small\ttfamily\color{axbluedk}[#1]}}
\newcommand{\ev}[1]{\langle #1\rangle}
\newcommand{\ket}[1]{|#1\rangle}
\newcommand{\Oid}{O_{\mathrm{id}}}
\newcommand{\Var}{\operatorname{Var}}
\newcommand{\Cov}{\operatorname{cov}}
\newcommand{\E}{\mathbb{E}}
\newcommand{\Tr}{\operatorname{Tr}}
\newcommand{\PRX}{\mathrm{PRX}}
\newcommand{\CZ}{\mathrm{CZ}}
\newcommand{\intu}{\paragraph{Intuicja.}}

\title{\bfseries Gdzie mitygacja błędów przestaje się opłacać\\[4pt]
\large\normalfont Matematyka, intuicja i metodologia --- w kolejności slajdów}
\author{Karina Leśkiewicz\\\small KN Axion, Politechnika Wrocławska}
\date{\today}

\begin{document}
\maketitle

\noindent Tekst idzie w kolejności prezentacji; w nagłówkach podsekcji na niebiesko jest
identyfikator slajdu, do którego się odnoszą. Pola \PH{} to liczby z okna \texttt{w01}, takie
same jak żółte pola na slajdach.

\tableofcontents
\clearpage

% =====================================================================
\section*{Notacja}
\addcontentsline{toc}{section}{Notacja}

$O=\bar Z=\frac15\sum_{q=1}^5Z_q$ to obserwabla (średnie $Z$ po pięciu kubitach), $\Oid$ jej
wartość bezszumowa, $E$ wartość zmierzona na sprzęcie, $\hat O$ wartość po mitygacji. Błąd na
punkcie to $\varepsilon=|\hat O-\Oid|$. $S$ to liczba shotów na obwód ($8192$), $M$ liczba ramek
twirlingu lub kopii ($6$), $n_{\text{train}}$ liczba obwodów treningowych CDR ($20$).

Każdy estymator ma błąd średniokwadratowy
\[
  \mathrm{MSE}=\E\big[(\hat O-\Oid)^2\big]=B^2+\Var\hat O,\qquad B=\E\hat O-\Oid .
\]
Szum sprzętu tworzy obciążenie $B$, skończona liczba shotów tworzy wariancję. Mitygacja
zmniejsza $B$ kosztem wariancji. Cała praca jest o tym, kiedy ta wymiana się opłaca.

Jeden shot daje $\bar z=\frac15\sum_qz_q$, $z_q\in\{\pm1\}$, z wariancją $V\le1$; dla
niezależnych kubitów $V=\frac1{25}\sum_q(1-\ev{Z_q}^2)\le0{,}2$. Po $S$ shotach błąd
standardowy to $\sqrt{V/S}$, czyli przy $S=8192$ około $0{,}005$ (najwyżej $0{,}011$).

% =====================================================================
\section{QEM, a nie QEC}

\subsection{Czym jest mitygacja \slid{qem\_what}}

Procesor kwantowy nie zwraca $\ev{O}$, tylko ciągi bitów; $\ev{O}$ estymujemy z wielu shotów.
Zaszumiony procesor realizuje zamiast unitarnej $U_C$ kanał $\tilde{\mathcal C}$, więc zamiast
$\Oid=\Tr[O\,U_C\rho_0U_C^\dagger]$ estymujemy $E=\Tr[O\,\tilde{\mathcal C}(\rho_0)]$. Metoda
mitygacji to przepis: wybierz obwody $C_1,\dots,C_K$, zmierz ich $E_k$ i zwróć
$\hat O=g(E_1,\dots,E_K)$ dla klasycznej funkcji $g$. Stan kwantowy pozostaje zaszumiony;
poprawiamy tylko liczbę odczytaną na końcu.

\intu Szum ,,ściska'' wynik w stronę pewnej stałej. Informacja o $\Oid$ dalej w nim jest, tylko
pomniejszona. Mitygacja ją rozciąga z powrotem, a razem z nią rozciąga szum shotów.

\textbf{Dlaczego koszt rośnie wykładniczo.} Przy depolaryzacji z prawdopodobieństwem $p$ na
bramkę i $N$ bramkach sygnał maleje jak $E=(1-p)^N\Oid\approx e^{-pN}\Oid$. Odzyskanie $\Oid$ to
dzielenie przez $e^{-pN}$; szum shotów dzielimy tak samo, więc wariancja rośnie jak $e^{2pN}$, a
z nią liczba shotów potrzebnych do ustalonej precyzji. Takagi i~in.\ (2022) oraz Quek i~in.\
(2024) pokazali, że to cecha każdej metody QEM, nie tylko tej prostej: koszt w ogólności rośnie
wykładniczo z głębokością. Punkt, w którym mitygacja przestaje się opłacać, zawsze więc
istnieje. Dla skali: 57 CZ po $\sim1\%$ błędu to $pN\approx0{,}57$, sygnał ściśnięty do
$e^{-0{,}57}\approx0{,}57$, a wariancja po pełnej korekcie co najmniej $e^{1{,}14}\approx3$ razy
większa.

\subsection{Dlaczego nie korekcja \slid{why\_not\_qec}}

Korekcja koduje kubit logiczny w wielu fizycznych i na bieżąco poprawia błędy. Błąd logiczny
kodu o odległości $d$ skaluje się jak
\[
  p_L\approx A\,(p/p_{\text{th}})^{\lfloor(d+1)/2\rfloor}.
\]
Poniżej progu $p_{\text{th}}$ większy kod wykładniczo tłumi błąd, powyżej wykładniczo go
wzmacnia. Kod powierzchniowy zużywa $\sim d^2$ kubitów fizycznych na logiczny, a każdy kod
wymaga pomiaru syndromów w trakcie obwodu i sprzężenia zwrotnego. Pięć kubitów nie wystarcza
nawet na najmniejszy kod korygujący dowolny błąd jednokubitowy ($[[5,1,3]]$ zużywa wszystkie
pięć na jeden kubit logiczny i nic nie zostaje na pomiar syndromów). QEM nie ma progu: działa przy
każdym szumie, dopóki zostaje sygnał, i potrzebuje tylko pomiaru na końcu. Płaci shotami zamiast
kubitami.

Wcześniej testowaliśmy kod detekcyjny $[[4,2,2]]$: wykrywa jeden błąd i odrzuca shoty z
niezerowym syndromem. Koszt efektywny to $1/\text{akceptacja}$ i rośnie z głębokością, więc
jakościowo nie zmienia obrazu.

% =====================================================================
\section{Metody}

\subsection{Trzy narzędzia \slid{methods}}

\textbf{ZNE} nie wie nic o idealnym wyniku: mierzy, jak wynik zmienia się przy \emph{większym}
szumie, i cofa się po tej krzywej do zera. \textbf{CDR} korzysta z obwodów prawie takich
samych jak nasz, których ideał umiemy policzyć (Cliffordy), i uczy się z nich mapy
zaszumione$\to$idealne. \textbf{Twirling Pauliego} sam niczego nie mityguje: nie zmienia
średniego błędu, zmienia jego kształt (koherentny na stochastyczny), tak żeby pasował do modelu,
który zakłada CDR.

\subsection{Model szumu: szum ściska sygnał \slid{noise}}

Kanał depolaryzujący na $n$ kubitach to $\mathcal D_p(\rho)=(1-p)\rho+p\,I/2^n$. Dla
obserwabli $\Tr[O\mathcal D_p(\rho)]=(1-p)\ev{O}+p f$ z $f=\Tr O/2^n$. Globalna depolaryzacja
komutuje z każdą unitarną, więc $N$ takich kanałów rozsianych po obwodzie można zebrać na
końcu:
\[
  E=f+(1-p)^N(\Oid-f)\approx f+(\Oid-f)e^{-pN}.
\]
Jeśli szum wzmocnimy $\lambda$ razy ($N\to\lambda N$) i oznaczymy $\gamma=pN$:
\[
  \boxed{E(\lambda)=f+(\Oid-f)\,e^{-\gamma\lambda}}
\]
$f$ to \emph{podłoga}: wartość przy pełnym zaszumieniu. Dla średniego $Z$ w czystej
depolaryzacji $f=0$, ale relaksacja $T_1$ ciągnie kubity do $\ket0$, czyli $Z=+1$, więc na
transmonach $f>0$. Dlatego $f$ mierzymy, a nie zakładamy. Obciążenie surowego pomiaru:
\[
  B=E(1)-\Oid=(\Oid-f)(e^{-\gamma}-1)\approx-\gamma(\Oid-f).
\]
Model jest dokładny tylko dla globalnej depolaryzacji. Szum lokalny, korelacje i błędy
koherentne go psują: dla ZNE krzywa nie jest dokładnie wykładnicza, dla CDR mapa nie jest
dokładnie liniowa. To źródło obu granic, o których mówimy dalej.

\begin{center}
\begin{tikzpicture}[x=1.3cm,y=2.6cm]
  \draw[-{Latex},thick] (0,0)--(5.6,0) node[right]{$\lambda$};
  \draw[-{Latex},thick] (0,0)--(0,1.15) node[above]{$E(\lambda)$};
  \draw[dashed,axgolddk] (0,0.15)--(5.3,0.15) node[right]{\small $f$};
  \draw[very thick,axbluedk,domain=0:5.2,samples=60] plot(\x,{0.15+0.8*exp(-0.45*\x)});
  \foreach \l in {1,3,5} \fill[axblue,draw=axink] (\l,{0.15+0.8*exp(-0.45*\l)}) circle (2.5pt);
  \fill[axgold,draw=axink] (0,0.95) circle (3pt) node[left]{\small $\Oid$};
  \foreach \l in {1,3,5} \node[below] at (\l,0){\small\l};
\end{tikzpicture}
\end{center}

\subsection{ZNE: folding i Richardson \slid{zne\_fold}}

\textbf{Folding.} $U\to U(U^\dagger U)^k$. Bez szumu to wciąż $U$, ale bramek jest
$\lambda=2k+1$ razy więcej, więc szum rośnie $\lambda$ razy. Używamy $\lambda\in\{1,3,5\}$.
Składamy obwód \emph{już skompilowany} do bramek natywnych i nie kompilujemy go ponownie, bo
optymalizator skróciłby $U^\dagger U$ albo przetrasował go inaczej i skala przestałaby być skalą.
Kopie rozdzielamy barierami.

\textbf{Ekstrapolacja.} Przez punkty $(\lambda_k,E_k)$ prowadzimy wielomian Lagrange'a i bierzemy
jego wartość w zerze:
\[
  \hat E(0)=\sum_kc_kE(\lambda_k),\qquad c_k=\prod_{j\ne k}\frac{\lambda_j}{\lambda_j-\lambda_k}.
\]
Dla $\{1,3,5\}$: $c_1=\frac{3\cdot5}{2\cdot4}=\frac{15}8$, $c_3=\frac{1\cdot5}{(-2)\cdot2}=-\frac54$,
$c_5=\frac{1\cdot3}{(-4)(-2)}=\frac38$. Zawsze $\sum c_k=1$ (stała jest odtwarzana dokładnie), więc
$f$ przechodzi przez ekstrapolację nietknięte.

\textbf{Ile obciążenia zostaje.} Błąd interpolacji Lagrange'a w zerze:
\[
  E(0)-\hat E(0)=\frac{E'''(\xi)}{3!}\prod_k(0-\lambda_k)=-\tfrac{15}{6}E'''(\xi),\qquad\xi\in(0,5).
\]
Dla $E=f+(\Oid-f)e^{-\gamma\lambda}$ mamy $|E'''|\le\gamma^3|\Oid-f|$, więc
\[
  |B_{\text{ZNE}}|\lesssim2{,}5\,\gamma^3|\Oid-f|\quad\text{wobec}\quad|B|\approx\gamma|\Oid-f| .
\]
Przy małym $\gamma$ ZNE zbija obciążenie z rzędu $\gamma$ do rzędu $\gamma^3$. Przy dużym
wielomian przestaje przybliżać wykładniczy zanik i zysk znika.

\subsection{Dolna granica ZNE: stała $\kappa$ \slid{zne\_kappa}}

\intu Ekstrapolacja to dźwignia: współczynniki są duże i mają różne znaki, więc mały szum shotów
w punktach pomiarowych zamienia się w duży szum w $\lambda=0$. Przy małej liczbie shotów ZNE
szkodzi, przy dużej pomaga.

\textbf{Podział równy.} Budżet $S$ dzielimy po równo na $|\Lambda|$ skal. Pomiary są niezależne,
więc
\[
  \Var\hat E_{\text{ZNE}}=\sum_kc_k^2\frac{V}{S/|\Lambda|}=|\Lambda|\Big(\sum_kc_k^2\Big)\frac VS,
  \qquad \Var E_{\text{raw}}=\frac VS,
\]
\[
  \mathrm{MSE}_{\text{ZNE}}-\mathrm{MSE}_{\text{raw}}=\kappa\frac VS+\big(B_{\text{ZNE}}^2-B^2\big),
  \qquad \kappa=|\Lambda|\sum_kc_k^2-1 .
\]
$\kappa$ mówi, ile razy więcej shotów ZNE potrzebuje do tej samej wariancji co pomiar surowy.
ZNE pomaga, gdy różnica jest ujemna:
\[
  S>S^*=\frac{\kappa V}{B^2-B_{\text{ZNE}}^2}.
\]
Im mniejsze surowe obciążenie (płytszy obwód), tym większe $S^*$. Dla skali: $V=0{,}2$,
$B=0{,}1$, $B_{\text{ZNE}}\approx0$ daje $S^*\approx290$; przy $B=0{,}02$ już $S^*\approx7300$.

\textbf{Podział Neymana.} Dajemy $n_k$ shotów na skalę $k$, $\sum n_k=S$, i minimalizujemy
$\sum_kc_k^2\sigma_k^2/n_k$. Mnożnik Lagrange'a daje $c_k^2\sigma_k^2/n_k^2=\mu$, czyli
$n_k\propto|c_k|\sigma_k$, i po podstawieniu
\[
  \Var_{\text{opt}}=\frac{(\sum_k|c_k|\sigma_k)^2}{S},\qquad
  \kappa_{\text{opt}}=\Big(\frac{\sum_k|c_k|\sigma_k}{\sigma_1}\Big)^2-1 .
\]
Więcej shotów idzie tam, gdzie dźwignia jest największa.

\begin{center}
\begin{tabular}{@{}lcccc@{}}
\toprule
$\Lambda$ & $c_k$ & $\sum c_k^2$ & $\kappa$ (równy) & $\kappa_{\text{opt}}$ (równe $\sigma$) \\
\midrule
$\{1,3\}$ & $\tfrac32,-\tfrac12$ & $2{,}5$ & $4{,}00$ & $3{,}0$ \\
$\{1,3,5\}$ & $\tfrac{15}8,-\tfrac54,\tfrac38$ & $5{,}22$ & $14{,}66$ & $11{,}25$ \\
$\{1,3,5,7\}$ & $\tfrac{35}{16},-\tfrac{35}{16},\tfrac{21}{16},-\tfrac5{16}$ & $11{,}39$ & $44{,}56$ & $35{,}0$ \\
\bottomrule
\end{tabular}
\end{center}

Każdy dodatkowy węzeł obniża obciążenie o rząd, ale mniej więcej potraja $\kappa$. Kluczowe:
$\kappa$ zależy tylko od węzłów i podziału shotów, nie od obwodu. To przewidywanie da się
sfalsyfikować. Mierzymy je jako wyraz wolny regresji
$(\mathrm{MSE}_{\text{ZNE}}-\mathrm{MSE}_{\text{raw}})\cdot S$ względem $S$ (podzielony przez $V$).

\subsection{Górna granica ZNE: zanik sygnału \slid{zne\_collapse}}

Gdy obwód jest tak głęboki, że $E(3)=E(5)=f$, to z $\sum c_k=1$:
\[
  \hat E(0)=c_1E_1+(1-c_1)f=\tfrac{15}8E_1-\tfrac78f .
\]
To nie jest ekstrapolacja, tylko przeskalowanie jednego pomiaru. Dla $f=0$ to $1{,}875\,E_1$;
gdy ściśnięcie przy $\lambda=1$ wynosiło akurat około $0{,}53$, wynik wygląda na świetną
mitygację, choć szum nie został poprawnie wzmocniony. Köster i Mauerer (2026) pokazali to na
54-kubitowym procesorze IQM. Ze skalarnego $\hat E$ tego nie widać, więc mamy trzy kontrole.

\textbf{Garbage folding.} $U(V^\dagger U)^k$, gdzie $V$ ma ten sam szkielet CZ i tę samą liczbę
PRX co $U$, ale losowe kąty. Koszt identyczny, funkcja zniszczona, więc $E(\lambda>1)$ jest z
założenia śmieciem. Zysk, który przetrwa tę kontrolę, jest artefaktem. Odwrotna kontrola,
\emph{identity folding}: ten sam koszt, a funkcja zachowana (każda bramka 1q w połówce $U^\dagger$
sparowana ze swoją odwrotnością, CZ komutują i dają identyczność). Ten wariant powinien zachowywać
się jak zwykły folding.

\textbf{Ujemna masa.} Richardsona stosujemy do każdego z $2^5=32$ prawdopodobieństw stanów:
\[
  \hat P(s)=\sum_kc_kP_s(\lambda_k),\qquad W_{\text{neg}}=\sum_s\max\big(0,-\hat P(s)\big).
\]
Normalizacja zawsze się zgadza ($\sum c_k=1$), ale prawdziwy rozkład nie ma ujemnych wpisów. Duża
ujemna masa oznacza, że krzywe $P_s(\lambda)$ nie są gładkie tak, jak zakłada Richardson. Szum
shotów, wzmocniony przez $c_3=-\frac54$, daje trochę ujemnej masy nawet przy idealnym szumie, więc
próg $W_{\text{neg}}\le0{,}02$ wzięliśmy z symulacji zerowej (depolaryzacja, 8192 shoty;
99.\ percentyl $0{,}014$).

\textbf{SNR.} $\mathrm{SNR}=|E(\lambda_{\max})|/\sqrt{V/N}$. Poniżej $\sim2$ największa skala jest
nieodróżnialna od zera.

Razem: ZNE ma \emph{okno}. Od dołu ogranicza je wariancja ($S>S^*$), od góry zanik sygnału
(SNR, $W_{\text{neg}}$, garbage).

\subsection{CDR \slid{cdr}}

\textbf{Cliffordy i Gottesman--Knill.} Bramki Clifforda przeprowadzają operatory Pauliego na
operatory Pauliego. Obwód Clifforda na $\ket0^{\otimes n}$ daje stan stabilizatorowy, opisany
przez $n$ generatorów, i symuluje się go w czasie wielomianowym. Dla takiego stanu
$\ev{P}\in\{-1,0,1\}$ dla każdego Pauliego $P$.

$\PRX(\theta,\varphi)=e^{-i\frac\theta2(\cos\varphi X+\sin\varphi Y)}$ jest Cliffordem, gdy
$\varphi$ i $\theta$ są wielokrotnościami $\pi/2$: oś to $\pm X$ lub $\pm Y$, obrót o
wielokrotność ćwierćobrotu. CZ jest Cliffordem.

\textbf{Zbiór treningowy.} Na skompilowanym obwodzie celu zastępujemy kąty każdej PRX wartościami
z $\{0,\frac\pi2,\pi,\frac{3\pi}2\}$. Obwód 0 dostaje kąty najbliższe; pozostałe losują je z wagą
$\propto e^{-d^2/\sigma^2}$, gdzie $d$ to odległość kątowa od oryginału, $\sigma=\pi/4$
(\texttt{near\_target}). Zbiór leży więc blisko celu i fit interpoluje zamiast ekstrapolować.
Szkielet CZ jest identyczny z celem z konstrukcji, więc obwody treningowe widzą ten sam szum CZ.
Ideał każdego liczymy symulatorem stabilizatorowym.

\textbf{Dlaczego prosta.} Dla depolaryzacji $E=(1-p)\Oid+pf$, więc
\[
  \Oid=aE+b,\qquad a=\frac1{1-p},\quad b=-\frac{pf}{1-p},
\]
\emph{ta sama} prosta dla każdego obwodu o tym samym szumie. Nauczona na Cliffordach przenosi się
na cel. Błąd odczytu też jest afiniczny: $\ev Z_{\text{zm}}=(1-p_{01}-p_{10})\ev Z+(p_{10}-p_{01})$,
więc CDR koryguje go przy okazji (dokładnie, gdy kubity mają podobne błędy odczytu).

\textbf{Fit.} Dla par $(x_i,y_i)=(E_i,O_{\mathrm{id},i})$:
\[
  a=\frac{\Cov(x,y)}{\Var x},\qquad b=\bar y-a\bar x,\qquad\hat O=aE_{\text{cel}}+b .
\]

\textbf{Gdzie to pęka.} Błąd koherentny nie ściska wyniku o stały czynnik. Jego efekt zależy od
interferencji w \emph{konkretnym} obwodzie: dla jednego błędy się dodają, dla innego znoszą.
Punkty $(x_i,y_i)$ rozjeżdżają się wokół prostej, $r^2$ spada, a prosta z Cliffordów gorzej
pasuje do celu.

\textbf{Rozcieńczenie regresji.} $x_i$ to średnie ze skończonej liczby shotów,
$x_i=x_i^\star+\eta_i$, $\Var\eta_i=\sigma_{x,i}^2$. Wtedy
$\E a_{\text{OLS}}\approx a\cdot\Var x^\star/(\Var x^\star+\bar\sigma_x^2)<a$: nachylenie jest
zaniżone i CDR niedokorygowuje. Poprawka errors-in-variables odejmuje znaną wariancję:
\[
  a_{\text{EIV}}=\frac{\Cov(x,y)}{\Var x-\bar\sigma_x^2}.
\]
Gdy prawdziwe $x$ prawie się nie różnią, mianownik dąży do zera i nachylenie wybucha; dlatego
poprawkę ograniczamy do $4\times$ względem OLS i raportujemy oba fity.

\subsection{Twirling Pauliego \slid{twirl}}

\intu Błąd koherentny to systematyczne przekręcenie o mały kąt, zawsze w tę samą stronę: przy
kolejnych bramkach dodają się amplitudy. Twirling przed każdą bramką losowo zmienia układ
odniesienia, a za nią go przywraca. Bramka robi to samo, ale błąd raz wypada w jedną stronę, raz
w drugą. Po uśrednieniu zostaje losowe rozmycie, w którym dodają się prawdopodobieństwa, nie
amplitudy.

\textbf{Funkcja bez zmian.} Dla Pauliego $A$ ($A^2=I$) i Clifforda CZ operator
$B=\CZ\,A\,\CZ^\dagger$ jest Paulim i
\[
  B\,\CZ\,A=\CZ\,A\,\CZ^\dagger\CZ\,A=\CZ\,A^2=\CZ .
\]
Np.\ $A=X\otimes I$ daje $B=X\otimes Z$.

\textbf{Szum po uśrednieniu.} Zaszumiona bramka to $\Lambda\circ\mathcal{CZ}$. Zakładamy bezbłędne
ramki (ich błąd $\sim0{,}1\%$ jest dziesięć razy mniejszy niż CZ). Gdy $A$ przebiega 16
Pauliego, $B$ też przebiega wszystkie, więc
\[
  \overline{\mathcal B\circ\Lambda\circ\mathcal{CZ}\circ\mathcal A}
  =\Big(\tfrac1{16}\textstyle\sum_B\mathcal B\circ\Lambda\circ\mathcal B\Big)\circ\mathcal{CZ}
  =\Lambda_{\text{tw}}\circ\mathcal{CZ}.
\]
Zapiszmy $\Lambda(\rho)=\sum_{\alpha\beta}\chi_{\alpha\beta}P_\alpha\rho P_\beta$. Diagonala
$\chi_{\alpha\alpha}$ to prawdopodobieństwa błędów Pauliego, elementy pozadiagonalne to
interferencja ścieżek błędu, czyli część koherentna. Ponieważ
$P_kP_\alpha P_k=(-1)^{\langle k,\alpha\rangle}P_\alpha$,
\[
  \Lambda_{\text{tw}}(\rho)=\sum_{\alpha\beta}\Big[\tfrac1{4^n}\textstyle\sum_k(-1)^{\langle k,\alpha\rangle+\langle k,\beta\rangle}\Big]\chi_{\alpha\beta}P_\alpha\rho P_\beta=\sum_\alpha\chi_{\alpha\alpha}P_\alpha\rho P_\alpha ,
\]
bo średnia znaków to $1$ dla $\alpha=\beta$ i $0$ w przeciwnym razie (połowa Pauliego komutuje z
$P_\alpha P_\beta$, połowa antykomutuje). Twirling zachowuje diagonalę, a więc i wierność, i
\emph{zeruje} część koherentną.

\textbf{Przykład $ZZ$.} $U=e^{-i\frac\theta2ZZ}=\cos\frac\theta2I-i\sin\frac\theta2ZZ$. Po
twirlingu: $\Lambda_{\text{tw}}(\rho)=\cos^2\frac\theta2\,\rho+\sin^2\frac\theta2\,ZZ\rho ZZ$.
$N$ powtórzeń błędu koherentnego to obrót o $N\theta$, niewierność
$\sin^2\frac{N\theta}2\approx N^2\theta^2/4$; po twirlingu $\approx N\theta^2/4$. Kwadratowo kontra
liniowo.

\textbf{W praktyce} losujemy $M=6$ ramek (niezależny Pauli dla każdego CZ) i uśredniamy. Każda ramka
to osobny obwód, więc koszt pomiarowy rośnie $M$ razy (stąd równy budżet).

\textbf{Dlaczego to pomaga CDR.} CDR zakłada jedną liniową mapę szumu dla wszystkich obwodów o tym
samym szkielecie. Szum koherentny zależy od obwodu i łamie to założenie; po twirlingu szum jest
kanałem Pauliego, którego wpływ jest bliższy wspólnemu ściśnięciu. Ostrożnie: twirling działa
empirycznie, ale mechanizm wykracza poza prosty błąd $ZZ$ na CZ. Model z samym koherentnym $ZZ$
nie odtwarza tego, co widzimy na szczeblach z routingiem.

\subsection{Logiczny vs natywny \slid{twirl\_kinds}}

Różnica to kolejność kroków. \emph{Logiczny}: ramki wokół bramek splątujących obwodu
logicznego, potem kompilacja. CZ, które kompilator dokłada w SWAP-ach, zostają bez ramki.
\emph{Natywny}: najpierw kompilacja, potem ramka wokół każdego fizycznego CZ.

\begin{center}
\begin{tabular}{@{}lcccc@{}}
\toprule
wzorzec & logiczne bramki 2q & CZ po routingu & CZ ze SWAP-ów & owinięte w logicznym \\
\midrule
star & 12 & 12 & 0 & 100\% \\
linear & 12 & 30 & 18 & 40\% \\
circular & 15 & 39 & 24 & 38\% \\
all-to-all & 30 & 57 & 27 & 53\% \\
\bottomrule
\end{tabular}
\end{center}

Na szczeblach z routingiem twirling logiczny obejmuje tylko 40--50\% CZ, więc jego zysk jest
dolnym ograniczeniem pełnego twirlingu. W wariancie logicznym Pauli to zwykłe bramki 1q przed
kompilacją, więc optymalizator łączy je z sąsiednimi obrotami. Porównanie obu (H3) robimy na
wspólnym zbiorze treningowym, żeby różnica wynikała tylko z miejsca ramek.

\subsection{Scalanie bramek 1q \slid{merge}}

Ramki to też bramki. Wstawione jako osobne impulsy wydłużają obwód: przy 57 CZ z 71 do 202 PRX.
Wtedy testujemy ,,twirling plus dodatkowy szum'' i stąd wziął się nasz pierwszy, ujemny wynik
natywnego twirlingu.

Naprawa korzysta z trzech faktów.
\begin{enumerate}
  \item Każde $U\in SU(2)$ to $R_Z(\alpha)\,\PRX(\theta,\varphi)$ z dokładnością do fazy (rozkład
  Eulera i $\PRX(\theta,\varphi)=R_Z(\varphi)R_X(\theta)R_Z(-\varphi)$). Ciąg dowolnie wielu
  bramek 1q to więc jeden PRX i jeden obrót $Z$.
  \item Obrót $Z$ nie wymaga impulsu: $R_Z(\alpha)\PRX(\theta,\varphi)R_Z(-\alpha)=\PRX(\theta,\varphi+\alpha)$,
  więc wchłaniamy go w fazę następnych PRX (,,wirtualne $Z$''). CZ i $R_Z$ są diagonalne i
  komutują, więc $R_Z$ przechodzi przez CZ do następnego ciągu.
  \item Pomiar jest w bazie $Z$, a $R_Z$ zmienia tylko fazy, więc końcowy obrót $Z$ odrzucamy.
\end{enumerate}
Po scaleniu między dwoma CZ na każdym kubicie jest co najwyżej jeden PRX (sprawdzane w kodzie:
najdłuższy ciąg bramek 1q ma długość 1; górne ograniczenie to $2N_{CZ}+n$ PRX). Szkielet CZ się
nie zmienia. Narzut twirlingu spada z $2{,}4$--$2{,}8\times$ do $1{,}23$--$1{,}35\times$ (przy 57
CZ: 92--95 PRX). Scalamy we \emph{wszystkich} ramionach, żeby liczby bramek były porównywalne.

\subsection{Równy budżet \slid{budget}}

Ramię z twirlingiem mierzy każdy obwód w $M$ ramkach. Porównane z CDR mierzonym raz wygrywa samym
uśrednianiem. Dlatego każde ramię wydaje na punkt ewaluacyjny
\[
  B=(n_{\text{train}}+1)\,M\,S=21\cdot6\cdot8192=1\,032\,192\ \text{shotów},
\]
a różnią się tylko tym, jak ten budżet dzielą (cel zawsze $M$ razy):

\begin{center}
\begin{tabular}{@{}llcl@{}}
\toprule
ramię & obwody treningowe & twirling & rola \\
\midrule
\texttt{off\_copies} & $n_{\text{train}}$ obwodów $\times$ $M$ identycznych kopii & nie & kontrola \\
\texttt{twirl\_logical} & $n_{\text{train}}\times M$ ramek, ramki przed kompilacją & tak & wariant \\
\texttt{twirl\_native} & $n_{\text{train}}\times M$ ramek, po kompilacji, scalone & tak & główny test \\
\texttt{spread\_off} & $n_{\text{train}}M$ różnych obwodów $\times$ 1 kopia & nie & kontrola \\
\texttt{spread} & $n_{\text{train}}M$ różnych obwodów $\times$ 1 ramka & tak & propozycja \\
\bottomrule
\end{tabular}
\end{center}

Średnia z $M$ pomiarów po $S$ shotów ma wariancję $V/(MS)$ niezależnie od tego, czy to kopie, czy
ramki. \texttt{off\_copies} kupuje więc tę samą redukcję wariancji co twirling, tylko bez
twirlingu, a różnica \texttt{twirl\_native}$-$\texttt{off\_copies} to czysty efekt twirlingu.
(Ramki dodają jeszcze wariancję między ramkami, więc porównanie jest raczej na niekorzyść
twirlingu.)

\subsection{Spread \slid{spread}}

Ten sam budżet można wydać na 20 obwodów po 6 pomiarów albo na 120 różnych po jednym. Przy
głębokich obwodach przeważa drugie.

\textbf{Dlaczego ideały są zwykle zerowe.} Dla stanu stabilizatorowego $\ev{Z_q}=\pm1$ tylko wtedy,
gdy $\pm Z_q$ należy do grupy stabilizatorów. Głęboki obwód Clifforda rozprasza stabilizatory na
operatory wielokubitowe, więc prawie zawsze $\ev{Z_q}=0$ i $\bar Z=0$. Przy 57 CZ dotyczy to
około 80\% obwodów treningowych.

\textbf{Degeneracja fitu.} Punkt z $y_i=0$ prawie nic nie mówi o nachyleniu. Jeśli niezerowy ideał
ma każdy obwód z prawdopodobieństwem $q\approx0{,}2$, liczba informatywnych punktów to
$\mathrm{Bin}(n,q)$. Dla $n=20$: średnio 4, a
$P(\le1)=0{,}8^{20}+20\cdot0{,}2\cdot0{,}8^{19}\approx0{,}07$, czyli w co czternastym punkcie
ewaluacyjnym prosta opiera się na najwyżej jednym niezerowym punkcie. Dla $n=120$: średnio 24,
$P(\le1)\sim10^{-10}$.

\textbf{Cena.} Punkt spread jest mierzony raz, więc jego wariancja to $V/S$ plus wariancja między
ramkami, zamiast $V/(MS)$. Rozcieńczenie regresji jest silniejsze i poprawka EIV ważniejsza.
Wariancji pojedynczego punktu nie da się zmierzyć, więc bierzemy model shotów
$\Var\bar Z\approx(1-\ev O^2)/(nS)$ plus wariancję między ramkami zmierzoną na celu.

% =====================================================================
\section{Obciążenie: model i obwód}

Ta sekcja rozwija slajdy \texttt{dataset}, \texttt{workload} i \texttt{odra}.

\subsection{Dane \slid{dataset}}

Zbiór \emph{two curves} z pakietu \texttt{qml\_benchmarks} (Bowles, Ahmed, Schuld 2024). Parametr
$t$ jest odwzorowany przez dwa wielomiany stopnia 8 w $\mathbb R^8$; druga krzywa jest przesunięta
o $0{,}1$, a do punktów dodany szum gaussowski $0{,}01$. Klasa $y\in\{-1,+1\}$ mówi, do której
krzywej należy punkt. Małe przesunięcie (wersja ,,trudna'') sprawia, że wiele punktów leży blisko
granicy decyzji, co jest potrzebne do pytania, czy mitygacja zmienia decyzje.

300 próbek, cechy przeskalowane do $[-\pi,\pi]$ (to kąty obrotów), podział warstwowy 180/120
(\texttt{test\_size}$=0{,}4$, ziarno 42). Z części testowej zamrażamy 100 punktów referencyjnych,
zbalansowanych klasowo.

\subsection{Ansatz: hardware-efficient z data re-uploadingiem \slid{workload}}

\textbf{Hardware-efficient ansatz} (Kandala i~in.\ 2017) to obwód zbudowany z warstw, które
łatwo wykonać na sprzęcie: obroty jednokubitowe i bramki dwukubitowe na sąsiednich kubitach, bez
struktury wynikającej z problemu. Jest płytki i ekspresywny przy małej liczbie bramek; ceną jest
brak indukcyjnego obciążenia i (dla głębokich, szerokich wersji) jałowe plateau gradientu. Przy 5
kubitach, 3 warstwach i obserwabli będącej sumą jednokubitowych $Z$ plateau nie jest problemem.

\textbf{Data re-uploading} (Pérez-Salinas i~in.\ 2020) to kodowanie danych \emph{w każdej warstwie},
a nie tylko na początku. Pojedyncze kodowanie $R_X(x)$ daje funkcję $x$ zawierającą tylko
częstotliwości $\{-1,0,1\}$, czyli $a+b\cos x+c\sin x$. Ogólnie (Schuld, Sweke, Meyer 2021)
model jest szeregiem Fouriera
\[
  f(x)=\sum_{\omega\in\Omega}c_\omega e^{i\omega x},
\]
gdzie widmo $\Omega$ tworzą różnice wartości własnych generatorów kodowania. Kodowanie
$R_X(w_jx)=e^{-iw_jxX/2}$ ma wartości własne generatora $\pm w_j/2$, więc daje różnice
$\{-w_j,0,w_j\}$, a kilka kodowań tej samej cechy daje
$\Omega=\{\sum_jk_jw_j:\ k_j\in\{-1,0,1\}\}$. Więcej kodowań to bogatsze widmo; trenowane skale
$w_j$ pozwalają dostroić częstotliwości do danych. Współczynniki $c_\omega$ ustalają trenowane
obroty.

\textbf{Dokładna budowa.} $n=5$ kubitów, $L=3$ warstwy, $d=8$ cech. Warstwa $\ell$:
\begin{enumerate}
  \item \emph{kodowanie}: na kubicie $q$ bramka $R_X(w^{\text{enc}}_{\ell q}\,x_{k\bmod 8})$, gdzie
  $k=5\ell+q$ to bieżący indeks (cechy idą cyklicznie po kubitach i warstwach);
  \item \emph{trenowane obroty}: $R_Y(w_{\ell q})\,R_Z(w'_{\ell q})$ na każdym kubicie;
  \item \emph{blok splątania}: bramki CX na parach z wzorca.
\end{enumerate}
Na końcu jeszcze jedna warstwa $R_YR_Z$. 15 slotów kodowania na 8 cech: $x_1,\dots,x_7$ są
kodowane dwukrotnie, $x_8$ raz. Liczba wag:
\[
  \underbrace{Ln}_{\text{skale}}+\underbrace{2Ln}_{R_Y,R_Z}+\underbrace{2n}_{\text{końcowe}}=15+30+10=55,
\]
taka sama dla każdego wzorca. Pojemność modeli jest więc dopasowana, a różni je tylko splątanie.

\begin{center}
\begin{tikzpicture}[x=1cm,y=0.62cm,
  g/.style={draw=axink,thick,fill=white,minimum width=1.25cm,minimum height=0.48cm,font=\scriptsize,inner sep=1pt},
  e/.style={draw=axink,thick,fill=axblue!40,minimum width=1.25cm,minimum height=0.48cm,font=\scriptsize,inner sep=1pt}]
  \foreach \q in {0,...,4} {
    \draw[thick] (-0.6,-\q)--(8.2,-\q);
    \node[left,font=\small] at (-0.6,-\q) {$q_\q$};
    \node[e] at (0.5,-\q) {$R_X(w x_k)$};
    \node[g] at (2.1,-\q) {$R_Y(w)$};
    \node[g] at (3.6,-\q) {$R_Z(w)$};
  }
  \foreach \t/\x in {1/5.0,2/5.8,3/6.6,4/7.4} {
    \draw[thick] (\x,0)--(\x,-\t);
    \fill (\x,0) circle (2.5pt);
    \draw[thick,fill=white] (\x,-\t) circle (5pt);
    \draw[thick] (\x,-\t+0.28)--(\x,-\t-0.28);
    \draw[thick] (\x-0.17,-\t)--(\x+0.17,-\t);
  }
  \node[font=\small] at (3.8,1.0) {jedna warstwa (powtarzana $L=3$ razy); wzorzec star, hub $q_0$};
\end{tikzpicture}
\end{center}

\textbf{Wzorce splątania} (pary w jednym bloku): star $(0,q)$ dla $q=1..4$ (4 pary), linear
$(q,q{+}1)$ (4), circular linear plus $(4,0)$ (5), all-to-all wszystkie $\binom52=10$ par.

\textbf{Wyjście i decyzja.} $f(x)=\ev{\bar Z}\in[-1,1]$, decyzja $\hat y=\mathrm{sign}\,f(x)$.
Margines punktu to $|\Oid|$, czyli odległość idealnego wyjścia od granicy decyzji.

\textbf{Trening.} Bezszumowo (\texttt{StatevectorEstimator}), funkcja straty MSE między $f(x)$ a
etykietą $\pm1$, cały zbiór treningowy w każdym kroku, Adam z $\eta=0{,}08$, 60 epok, wagi
początkowe $0{,}1\cdot\mathcal N(0,1)$, ziarno 42. Gradient z reguły przesunięcia parametru: dla
bramki $e^{-i\theta P/2}$
\[
  \frac{\partial\ev O}{\partial\theta}=\frac12\Big[\ev O_{\theta+\pi/2}-\ev O_{\theta-\pi/2}\Big],
\]
dokładnie, nie jako różnica skończona; dla kodowania $R_X(wx)$ dochodzi czynnik $x$ z reguły
łańcuchowej. Każdy wzorzec ma osobny trening na tych samych danych; wagi są potem zamrożone.

\begin{center}
\begin{tabular}{@{}lcccc@{}}
\toprule
 & star & linear & circular & all-to-all \\
\midrule
logiczne bramki 2q & 12 & 12 & 15 & 30 \\
dokładność trening / test (bez szumu) & 0,778 / 0,725 & 0,744 / 0,692 & 0,733 / 0,675 & 0,772 / 0,767 \\
\bottomrule
\end{tabular}
\end{center}

To cztery modele o tej samej architekturze, a nie jeden model na czterech topologiach. Dzięki
temu każdy szczebel jest prawdziwym klasyfikatorem, a nie wagami gwiazdy przeniesionymi na inne
splątanie. 0,725 na slajdzie to model star.

\subsection{Bramki natywne i kompilacja \slid{odra}}

ODRA wykonuje $\PRX(\theta,\varphi)$ (impuls mikrofalowy o fazie $\varphi$, obrót o $\theta$ wokół
osi w płaszczyźnie $xy$) i CZ. Obroty $Z$ realizujemy wirtualnie, zmieniając fazę kolejnych
impulsów. Bramki ansatzu tłumaczą się tak:
\[
  R_X(\alpha)=\PRX(\alpha,0),\quad R_Y(\alpha)=\PRX(\alpha,\tfrac\pi2),\quad
  R_Z\ \text{wirtualne},\quad \mathrm{CX}=(I\otimes H)\,\CZ\,(I\otimes H),
\]
a $H$ to PRX i obrót $Z$. Po scaleniu każda warstwa daje na kubicie jeden PRX przed blokiem
splątania (dowolna unitarna 1q to jeden PRX), plus PRX z rozkładu CX.

Kompilacja ma jedną ścieżkę: \texttt{transpile} do bazy $\{\texttt{r},\texttt{cz}\}$ na mapie
gwiazdy, układ logiczny $=$ fizyczny (logiczny kubit 0 to hub QB3), \texttt{optimization\_level=1},
routing z ustalonym \texttt{seed\_transpiler=42}, potem scalanie PRX. Wynik jest deterministyczny
co do bitu. (Bez ustalonego ziarna ten sam obwód kompilował się za każdym razem do innego
szkieletu, a drabina wychodziła 12/27/36/54.)

\textbf{Routing na gwieździe.} CZ działa tylko na krawędziach z hubem. Bramka między dwoma liśćmi
wymaga SWAP, a $\mathrm{SWAP}=\mathrm{CX}_{12}\mathrm{CX}_{21}\mathrm{CX}_{12}$ to 3 CZ. Tak powstaje
drabina 12/30/39/57 CZ: ta sama architektura, rosnący koszt dwukubitowy.

\textbf{Parametry urządzenia.} Błąd CZ z interleaved RB to $0{,}7$--$1{,}2\%$, błąd PRX z RB
$\sim0{,}1\%$, czyli dziesięć razy mniej, więc oś badania to liczba CZ. $T_2^*$ wynosi 3--6\,$\mu$s.
Kubity czekające na hub tracą w tym czasie koherencję jak $e^{-t/T_2^*}$.

% =====================================================================
\section{Metodologia}

Pełny opis eksperymentu w formie, w jakiej trafi do artykułu. Odpowiada slajdom \texttt{problem}
i \texttt{methodology}.

\subsection{Pytanie i projekt \slid{problem}}

Dla metody $m$, budżetu $B$ i liczby CZ $N$ wyznaczamy $\varepsilon_m(B,N)$ i pytamy, przy jakim
$N$ zmienia się kolejność metod. Ramiona CDR porównujemy \emph{przy tym samym} $B$
($1\,032\,192$ shoty na punkt); ZNE ($24\,576$ shotów na punkt) i pomiar surowy mają inny koszt, który
raportujemy obok błędu, a nie wyrównujemy. Czynniki projektu:
\begin{itemize}
  \item \textbf{szczebel} (12, 30, 39, 57 CZ), wewnątrz obiektu: każdy punkt jest mierzony na
  każdym szczeblu;
  \item \textbf{metoda}: pięć ramion CDR (tabela w sekcji o budżecie), ZNE w trzech wariantach
  foldingu (genuine, identity, garbage) i pomiar surowy;
  \item \textbf{punkt ewaluacyjny}: czynnik blokujący; wszystkie porównania są parowane po
  (szczebel, punkt).
\end{itemize}

\subsection{Punkty ewaluacyjne}

Ze 100 zamrożonych punktów referencyjnych wybieramy 24, warstwowo po marginesie $|\Oid|$ modelu
star: cztery kwartyle, po 6 punktów równo rozmieszczonych w randze wewnątrz kwartyla. Wybór jest
deterministyczny i zapisany w \texttt{campaign/eval\_points.json}. Okno używa 8 z nich (zestaw 0:
indeksy 0, 3, 6, \dots, 21, po 2 z kwartyla), kolejne okna rotują zestawy. Warstwowanie
gwarantuje, że w każdym oknie są punkty łatwe i trudne. Wcześniejsze kampanie brały pierwsze 4
punkty i nie miały rozkładu błędu na punkt.

\subsection{Okno pomiarowe i noce \slid{methodology}}

\begin{center}
\begin{tabular}{@{}lll@{}}
\toprule
noc & szczeble & zawartość \\
\midrule
A & star (12), circular (39) & 5 ramion CDR, ZNE $\{1,3,5\}$ $\times$ 3 warianty foldingu \\
P5 & star z wstawkami, L6/L9/L13 & tylko \texttt{off\_copies}, bez ZNE \\
B & linear (30), all-to-all (57) & jak A \\
\bottomrule
\end{tabular}
\end{center}
Wszystkie ramiona danego szczebla są w tej samej nocy, czyli w tej samej kalibracji. Szczeble są
łączone w pary płytki + głęboki, żeby każda noc pokrywała zakres osi.

\textbf{Randomizacja bloków.} Blok to (szczebel, ramię). Kolejność bloków w nocy jest losowana
ziarnem okna i zapisana w manifeście, więc dryf urządzenia w ciągu nocy nie ustawia się
systematycznie przeciwko któremuś ramieniu.

\textbf{Kontrola kalibracji.} Przed i po każdym bloku odczytujemy identyfikator zestawu kalibracji
(\texttt{calibration\_set\_id}). Jeśli zmienia się w trakcie bloku albo względem pierwszego bloku
nocy, okno jest oznaczane jako \texttt{split}. Analiza główna zostawia blok, analiza wrażliwości
usuwa szczeble, których to dotyczy.

\subsection{Konfiguracja metod}

\textbf{Pomiar surowy}: cel w $M=6$ kopiach po $S=8192$ shotów. \textbf{ZNE}: skale $\{1,3,5\}$,
folding globalny na obwodzie natywnym, Richardson bez obcinania (wartości poza $[-1,1]$ są
raportowane jako niefizyczne, nie poprawiane), trzy warianty foldingu o identycznym koszcie.
\textbf{CDR}: $n_{\text{train}}=20$ (z kotwicą), próbnik \texttt{near\_target} albo
\texttt{shared\_random} (jeden losowy zbiór wspólny dla ramion, do porównania logiczny vs natywny),
fit OLS jako główny, EIV jako drugorzędny. Zbiór treningowy jest losowany ziarnem zależnym od
(ziarno bazowe, punkt, okno), więc każdy punkt ma własny zbiór.

\subsection{Estymacja}

Z każdego shotu liczymy $\bar z$; wartość obwodu to średnia po shotach, a dla ramion z kopiami lub
ramkami średnia po $M$ obwodach (zliczenia ramek łączymy). Brak osobnej korekty błędu odczytu: CDR
koryguje jego afiniczną część, a ZNE go nie usuwa (pomiar jest jeden, więc błąd odczytu przechodzi
przez Richardsona jak stała).

\textbf{Metryka główna} to $\varepsilon=|\hat O-\Oid|$ na punkcie. Dokładność (zgodność znaku) jest
drugorzędna, bo jest dyskretna i przy 8 punktach na okno bardzo zaszumiona.

\subsection{Analiza statystyczna}

\textbf{Pary i logarytm.} Porównując ramiona A i B, liczymy dla każdego (szczebel, punkt)
$d_i=\log\varepsilon_{A,i}-\log\varepsilon_{B,i}$; noce A i B dają $n=32$ pary. Logarytm, bo efekty
są multiplikatywne, a błędy skośne. Błąd równy 0 podłogujemy na $10^{-4}$. Wielkość efektu to
średnia geometryczna ilorazu $\exp(\bar d)$.

\textbf{Test Wilcoxona znakowanych rang.} Rangujemy $|d_i|$ i sumujemy rangi dodatnich różnic,
$W^+$. Przy $H_0$ (rozkład $d$ symetryczny wokół zera) znaki są losowe, więc
$\E W^+=n(n+1)/4=264$ i $\mathrm{sd}\,W^+=\sqrt{n(n+1)(2n+1)/24}\approx53{,}6$. Test jest
jednostronny, w kierunku hipotezy. Wilcoxon nie zakłada normalności, tylko symetrię różnic, co na
skali logarytmicznej jest rozsądne.

\textbf{Korekta Holma.} Rodzina pięciu testów głównych. Bez korekty szansa choćby jednego
fałszywego wyniku sięga $1-0{,}95^5\approx23\%$. Holm: sortujemy $p_{(1)}\le\dots\le p_{(5)}$ i
odrzucamy, dopóki $p_{(i)}\le0{,}05/(6-i)$ (progi $0{,}01$; $0{,}0125$; $0{,}0167$; $0{,}025$;
$0{,}05$). Kontroluje to błąd rodzinny na poziomie $5\%$ i jest mocniejsze niż Bonferroni.

\begin{center}
\small
\begin{tabular}{@{}lp{5.1cm}p{4.6cm}p{3.6cm}@{}}
\toprule
id & hipoteza & porównanie & kryterium \\
\midrule
H1 & twirling obniża błąd CDR przy równym budżecie & \texttt{twirl\_native} vs \texttt{off\_copies} (\texttt{near\_target}) & $p_{\text{Holm}}<0{,}05$ i iloraz $<1$ \\
H3 & natywny co najmniej tak dobry jak logiczny & \texttt{twirl\_native} vs \texttt{twirl\_logical} (\texttt{shared\_random}) & $p_{\text{Holm}}<0{,}05$; inaczej iloraz i 95\% CI bootstrap \\
H6 & spread lepszy niż CDR z twirlingiem & \texttt{spread} vs \texttt{twirl\_native} & $p_{\text{Holm}}<0{,}05$ i iloraz $<1$ \\
H7a & przy stałym czasie więcej CZ to większy błąd & regresja P5, współczynnik $n_{CZ}$ & $>0$, $p_{\text{Holm}}<0{,}05$ \\
H7b & przy stałych bramkach dłuższy czas to większy błąd & regresja P5, współczynnik czasu & $>0$, $p_{\text{Holm}}<0{,}05$ \\
\midrule
H4 & okno ZNE na szczeblu (dodatkowa) & ZNE vs raw i vs garbage; $W_{\text{neg}}$, SNR & wszystkie 4 warunki \\
H5 & zmiany decyzji tylko przy małym marginesie (dodatkowa) & McNemar, raw vs \texttt{twirl\_native} & zmiany poniżej mediany $|\Oid|$ \\
\bottomrule
\end{tabular}
\end{center}

H4 jest spełniona na szczeblu, gdy: ZNE ma mniejszy błąd niż raw (jednostronny Wilcoxon po 8
punktach, $p<0{,}05$ bez korekty), mniejszy niż garbage na tym samym szczeblu, mediana
$W_{\text{neg}}\le0{,}02$ i mediana SNR przy $\lambda=5$ $\ge2$. Przewidywanie zapisane z góry: ZNE
płaci przy 12 CZ, a przy 57 CZ nie. Stabilności między kalibracjami ani $\kappa$ w tym oknie nie
testujemy; wartości $\kappa$ pochodzą z pilota i są eksploracyjne.

\textbf{Wykluczenia.} Żadnego punktu nie wykluczamy ze względu na wartość. Szczebel bez któregoś
ramienia wypada tylko z porównania z tym ramieniem (raportujemy zmniejszone $n$). Obwody
powtórzone po przerwanym zadaniu mają te same ziarna i są używane tak, jak zapisano; niepełne
rekordy odkładamy. Jeśli noc się nie skończy, testujemy to, co zmierzono, i raportujemy braki.

\subsection{Pre-rejestracja}

Plan analizy (\texttt{preregistration.md}) i kod, który go wykonuje (\texttt{analysis/prereg\_w01.py}),
zostały zacommitowane i oznaczone tagiem \texttt{prereg-v1} przed policzeniem czegokolwiek z
\texttt{w01}. Zapisane są hipotezy, testy, kryteria sukcesu, wykluczenia i progi. Chroni to przed
dobieraniem testu, podzbioru albo metryki po zobaczeniu danych. Każde odstępstwo raportujemy jako
odstępstwo.

Status, uczciwie: to plan analizy, a nie rejestracja przed zebraniem danych, bo noc A już trwała,
gdy plan powstawał. Żaden błąd, iloraz, fit ani test z \texttt{w01} nie był wtedy liczony. Okna
pilotażowe były analizowane i z nich pochodzą hipotezy; są eksploracyjne i nie wchodzą do testów
potwierdzających.

\subsection{Odtwarzalność}

Każdy obwód ma deterministyczny identyfikator i własne ziarno (ramki, bliźniaki garbage, zbiory
treningowe). Zapisujemy surowe zliczenia każdego obwodu (\texttt{counts.jsonl}) oraz manifest z
commitem, kolejnością bloków i ID kalibracji. Dla obwodów z twirlingiem zapisujemy ramkę (listę
Pauliego dla każdego CZ), z której razem z obwodem natywnym da się odtworzyć dokładnie wykonany
obwód. Tabele \texttt{cdr.csv} i \texttt{zne.csv} są pochodne i odtwarzane ze zliczeń. Ten sam
obwód można puścić w symulatorze z modelem szumu z kalibracji.

\subsection{Noc P5: co napędza degradację}

Głębszy obwód zmienia naraz liczbę CZ, liczbę PRX i czas trwania. Rozdzielamy je na jednym modelu
(star, 12 CZ). \texttt{cz\_pairs}: wstawiamy $k\in\{6,12,20\}$ par $\CZ\cdot\CZ=I$ na krawędziach
istniejących CZ (24, 36, 52 CZ; funkcja i PRX bez zmian, czas dłuższy o $2k$ czasów CZ).
\texttt{idle\_matched}: w tych samych miejscach przerwa o długości dwóch CZ na obu kubitach
(bramki bez zmian, czas jak w \texttt{cz\_pairs}). Pary są otoczone barierami, żeby nic ich nie
skasowało, a miejsca wstawienia zależą tylko od indeksu CZ, więc trafiają w to samo miejsce w celu
i w obwodach treningowych. L6/L9/L13 dokładają PRX (to inne klasyfikatory, dlatego analiza główna
jest na samych L3). Kontrasty:
\[
  \texttt{cz\_pairs}-\texttt{idle\_matched}=\text{efekt CZ przy stałym czasie},\qquad
  \texttt{idle\_matched}-\texttt{base}=\text{efekt czasu przy stałych bramkach}.
\]
Model:
\[
  \log\varepsilon_{ij}=\alpha_i+\beta_{CZ}\frac{n_{CZ,j}}{10}+\beta_t\,t_j[\mu\mathrm s]+u_{ij},
\]
z efektami stałymi punktu $\alpha_i$ i błędami standardowymi klastrowanymi po punkcie. $e^{\beta}$
to mnożnik błędu na $+10$ CZ albo $+1\,\mu$s. Czas $t_j$ to długość z harmonogramu ASAP (bramki
równoległe liczone raz, oczekiwanie wliczone), bo to ten czas widzi relaksacja. Postać
log-liniowa wynika z modelu ściśnięcia: czynnik $e^{-\sum p}$ jest wykładniczy w liczbie bramek i
w czasie.

% =====================================================================
\section{Wyniki}

Stan danych: okno \texttt{w01}, noc A (12 i 39 CZ, po 8 punktów) i noc P5 są zmierzone w całości.
Noc B (30 i 57 CZ) jest w toku, więc testy pre-rejestrowane H1, H3, H6 (32 pary, Holm) jeszcze nie
są policzone. Liczby z nocy A to opis połowy danych, nie test. H7 (P5) jest już rozstrzygnięte.
Kalibracja w nocy A była stała; w P5 zmieniła się w trakcie bloku \texttt{star\_cz52} (okno
\texttt{split}; analiza główna go zostawia, analiza wrażliwości usunie).

\subsection{Twirling i CDR \slid{res\_twirl}}

\begin{center}
\begin{tabular}{@{}lcccc@{}}
\toprule
CZ & \texttt{off\_copies} & \texttt{twirl\_native} & wygrane pary & iloraz geom. \\
\midrule
12 & 0,022 & 0,015 & 4/8 & 1,41 \\
39 & 0,053 & 0,040 & 7/8 & 0,68 \\
30, 57 & \PH & \PH & \PH & \PH \\
\bottomrule
\end{tabular}
\end{center}

Przy 39 CZ twirling obniża błąd CDR o około 30\% przy identycznym budżecie (7 z 8 par, test znaków
$p=0{,}035$, Wilcoxon $p=0{,}07$). Przy 12 CZ efektu nie ma: średni błąd spada, ale iloraz
geometryczny rośnie. To spodziewane: błąd CDR przy 12 CZ to około $0{,}02$, a szum shotów celu
($\sqrt{V/(MS)}\approx0{,}002$) mnożony przez nachylenie fitu i szum samego fitu daje błąd tego samego
rzędu. Twirling usuwa błąd systematyczny, a tu zostaje głównie statystyczny. Razem na 16 parach:
11 wygranych, $p=0{,}35$ --- test rozstrzygnie noc B, gdzie są głębsze szczeble.

Logiczny vs natywny na wspólnym zbiorze (\texttt{shared\_random}): logiczny $0{,}019$ / $0{,}025$,
natywny $0{,}031$ / $0{,}032$ (12 / 39 CZ), natywny wygrał 7 z 16 par. Przewagi natywnego (H3) na
razie nie ma. Pilot (3 okna): 11 z 12 par, iloraz $0{,}51$ (logiczny) i $0{,}54$ (natywny).

\subsection{Która metoda kiedy \slid{res\_methods}}

\begin{center}
\begin{tabular}{@{}lcccccc@{}}
\toprule
CZ & raw & ZNE & ZNE garbage & CDR + twirl & spread \\
\midrule
12 & 0,074 & 0,045 & 0,137 & 0,015 & 0,012 \\
39 & 0,161 & 0,112 & 0,103 & 0,040 & 0,068 \\
shoty / punkt & 8\,192 & 24\,576 & 24\,576 & 1\,032\,192 & 1\,032\,192 \\
\bottomrule
\end{tabular}
\end{center}

\textbf{Ważne dla metodologii: koszt.} Równy budżet obowiązuje \emph{między ramionami CDR}
(wszystkie po $1\,032\,192$ shoty na punkt). ZNE $\{1,3,5\}$ zużywa $3\cdot8192=24\,576$ shotów na
punkt, czyli około 42 razy mniej niż CDR. Porównanie ZNE z CDR nie jest więc porównaniem przy
równym koszcie; koszt raportujemy obok błędu. Na slajdzie ,,problem'' tytuł brzmi teraz
,,Porównujmy metody, licząc ten koszt'', a nie ,,przy tym samym koszcie''.

\textbf{ZNE i H4.} Przy 12 CZ ZNE jest prawdziwe: genuine jest wyraźnie lepsze od garbage
($0{,}045$ vs $0{,}137$, $p=0{,}01$), mediana $W_{\text{neg}}=0{,}005$, SNR $=23$. Poprawa względem
raw ($0{,}069\to0{,}045$) jest jednak na 8 punktach nieistotna ($p=0{,}47$). Przy 39 CZ ZNE jest
istotnie lepsze od raw ($p=0{,}004$), ale garbage daje to samo ($0{,}103$ vs $0{,}112$). To dokładnie
degeneracja z sekcji o górnej granicy: przy zaniku sygnału Richardson ,,poprawia'' wynik bez względu
na to, co jest w obwodach $\lambda=3,5$. Ważna obserwacja: $W_{\text{neg}}=0$ i SNR $=4$ tego
\emph{nie wykryły}; wykryła to tylko kontrola garbage. Formalnie H4 nie jest spełniona na żadnym
z dwóch szczebli (przy 12 CZ brak warunku ZNE $<$ raw, przy 39 CZ brak warunku ZNE $<$ garbage).

\textbf{Test degeneracji (floor tracking).} Liczymy, jak daleko $\hat E$ leży od przewidywania
degeneracji $c_1E_1+(1-c_1)f$ przy zmierzonej podłodze $f$. Przy 39 CZ mediana odległości to
$0{,}016$ (retencja sygnału $0{,}24$), czyli ZNE praktycznie \emph{jest} przeskalowanym pierwszym
pomiarem. Przy 12 CZ odległość to $0{,}17$ (retencja $0{,}82$): tam ekstrapolacja faktycznie
korzysta z $\lambda=3,5$. To ilościowe potwierdzenie wniosku z kontroli garbage.

\textbf{Identity folding nie jest dobrą kontrolą.} Wariant identity ma tę samą liczbę CZ i PRX co
genuine na każdej skali (sprawdzone w zapisanych obwodach), a w symulacji bezszumowej daje dokładnie
ideał. Na sprzęcie $\ev O$ ucieka jednak w stronę $+1$: przy 12 CZ $0{,}16$ i $0{,}28$, przy 39 CZ
$0{,}40$ i $0{,}57$ dla $\lambda=3,5$ (genuine zostaje przy ideale bliskim $0$). Powód leży w
konstrukcji: połowa bloku to same CZ bez bramek 1q pomiędzy, a druga połowa to pary
bramka--odwrotność. Bramki diagonalne nie mieszają populacji, więc błędy fazowe są tam niewidoczne w
pomiarze $Z$, a relaksacja $T_1$ (ciągnąca do $\ket0$, czyli $Z=+1$) działa bez przeszkód. Blok
identity ma więc ten sam koszt w bramkach, ale \emph{inny} szum niż $U^\dagger U$. Nie używamy go
do wniosków; kontrolą, która działa, jest garbage.

\textbf{CDR i spread.} CDR z twirlingiem ma najmniejszy błąd na obu szczeblach (przy wspomnianym
koszcie). Spread jest najlepszy przy 12 CZ, ale przy 39 CZ gorszy od twirlingu ($0{,}068$ vs
$0{,}040$, mediana $r^2$ $0{,}69$ wobec $0{,}92$). Przewidywana przewaga spread dotyczy głównie
najgłębszych szczebli (degeneracja fitu przy 57 CZ), więc H6 rozstrzygnie noc B.

\subsection{Co napędza degradację \slid{res\_depth}}

\begin{center}
\begin{tabular}{@{}llccc@{}}
\toprule
wynik & czynnik & zmiana błędu & 95\% CI & $p$ \\
\midrule
surowy & $+10$ CZ (stały czas) & $\times0{,}96$ & $0{,}77$--$1{,}19$ & $0{,}68$ \\
surowy & $+1\,\mu$s (stałe bramki) & $\times1{,}43$ & $1{,}30$--$1{,}58$ & $<0{,}0001$ \\
surowy & $+10$ PRX (z L6--L13) & $\times1{,}12$ & $0{,}81$--$1{,}56$ & $0{,}50$ \\
po CDR & $+10$ CZ & $\times1{,}34$ & $1{,}11$--$1{,}61$ & $0{,}002$ \\
po CDR & $+1\,\mu$s & $\times1{,}32$ & $0{,}77$--$2{,}27$ & $0{,}31$ \\
\bottomrule
\end{tabular}
\end{center}

H7b potwierdzona: przy stałych bramkach każda mikrosekunda podnosi surowy błąd o około 43\%. Przy
$p<0{,}0001$ przejdzie każdy próg Holma. H7a niepotwierdzona ($p=0{,}68$). W parach punkt po punkcie
dodatkowe CZ podnoszą surowy błąd w 6 z 8 punktów na każdym poziomie (zwykle o 20--50\%), ale 1--2
punkty z przypadkowo bliskim zera błędem surowym (np.\ iloraz $0{,}03$) ściągają log-regresję. Uczciwe
sformułowanie: wpływ CZ przy stałym czasie jest mały i niejednoznaczny, wpływ czasu duży i
jednoznaczny.

\textbf{Analiza wrażliwości (pre-rejestracja).} Bez bloku \texttt{star\_cz52}, w którym zmieniła
się kalibracja: surowy błąd $+1\,\mu$s $\times1{,}48$ ($p<0{,}0001$), $+10$ CZ $\times1{,}03$
($p=0{,}89$); po CDR $+10$ CZ $\times1{,}59$ ($p=0{,}0001$). Wnioski się nie zmieniają.

\textbf{Fit EIV (drugorzędny).} Prawie wszędzie daje to samo co OLS (różnice $\le0{,}007$). Wyjątek:
spread przy 39 CZ, gdzie EIV pogarsza błąd ($0{,}098$ wobec $0{,}068$). Mała prawdziwa rozpiętość $x$
przy dużym szumie pojedynczego pomiaru to dokładnie sytuacja, w której poprawka EIV jest niestabilna
(stąd limit $4\times$). OLS zostaje fitem głównym, tak jak w planie.

Po CDR obraz się odwraca. Zanik w czasie (relaksacja, defazowanie) zachowuje się jak globalne
ściśnięcie sygnału, które CDR koryguje prostą. Zostaje część od CZ, której CDR nie usuwa. To spójne z
tym, że twirling (który działa na szum CZ) pomaga CDR przy głębszych szczeblach. Modele L6/L9/L13
wypadają gorzej niż wstawki CZ o tej samej liczbie CZ i podobnym czasie (po CDR przy 36 CZ: L9
$0{,}135$, \texttt{cz36} $0{,}072$), co potwierdza wniosek z pilota.

\subsection{Stała $\kappa$ \slid{res\_kappa}}

Pilot ODRA (star, 12 CZ). ,,Przewidziane'' dla podziału równego to czysta liczba z węzłów; dla
Neymana to $\kappa_{\text{opt}}$ ze zmierzonymi $\sigma_k$. Stąd $3{,}06$ i $11{,}6$ zamiast $3{,}0$
i $11{,}25$: na większych skalach $\ev O$ jest bliżej zera, więc wariancja shotu jest inna.

\begin{center}
\begin{tabular}{@{}llccc@{}}
\toprule
węzły & podział & zmierzone & przewidziane & różnica \\
\midrule
$\{1,3,5\}$ & równy & $15{,}8\pm1{,}7$ & $14{,}66$ & $0{,}7\sigma$ \\
$\{1,3,5\}$ & Neyman & $12{,}5\pm1{,}1$ & $11{,}6$ & $0{,}8\sigma$ \\
$\{1,3\}$ & Neyman & $1{,}3\pm1{,}3$ & $3{,}06$ & $1{,}4\sigma$ \\
$\{1,3,5,7\}$ & Neyman & $26{,}0\pm3{,}3$; $25{,}7\pm11{,}3$ & $\sim37$ & $>3\sigma$ (pierwsze okno) \\
\bottomrule
\end{tabular}
\end{center}

Dla $\{1,3,5\}$ teoria się zgadza, a Neyman obniża $\kappa$ o około 20\%; o tyle samo spada próg
$S^*$. $\{1,3\}$ ma zbyt duży rozrzut na wniosek. $\{1,3,5,7\}$ jest wyraźnie poniżej
przewidywania; możliwe przyczyny to model wariancji przy $\lambda=7$, korelacje między skalami w
jednym oknie albo nieliniowość zależności MSE od $S$. To pytanie otwarte.

\subsection{Decyzje klasyfikatora \slid{res\_decisions}}

Decyzja to znak $\ev O$. Zmienia się z $E$ na $\hat O$ tylko wtedy, gdy $|\hat O-E|\ge|E|$; korekta
jest ograniczona, więc zmiany są możliwe tylko przy małych $|E|$, czyli małych marginesach.

\textbf{McNemar.} Każdy punkt ma dwie decyzje, każda zgodna lub nie z decyzją idealnego
klasyfikatora. Liczą się tylko pary niezgodne: $b$ (mitygacja zepsuła) i $c$ (naprawiła). Przy $H_0$
każda niezgodność to rzut monetą, $b\sim\mathrm{Bin}(b+c,\frac12)$:
\[
  p=\min\Big(1,\ 2\sum_{j=0}^{\min(b,c)}\binom{b+c}j2^{-(b+c)}\Big).
\]
Pilot (100 punktów na szczebel), CDR: 12 CZ $4/1$, $p=0{,}375$; 30 CZ $0/5$, $p=0{,}0625$; 39 CZ
$0/10$, $p=0{,}002$; 57 CZ $2/9$, $p=0{,}065$ (bez korekty). Wszystkie 31 zmian przy CDR i 18 przy
ZNE wypadło w połowie punktów o marginesie poniżej mediany; przy niezależności od marginesu szansa
byłaby rzędu $2^{-31}$.

\textbf{w01 noc A:} 16 punktów, zgodność raw $15/16$, po CDR $16/16$, po ZNE $16/16$. Jedyna zmiana
decyzji (naprawa) dotyczy punktu z marginesem $0{,}028$, przy medianie marginesu $0{,}16$--$0{,}31$.
Ten sam wzór co w pilocie.

\section{Wnioski i ograniczenia}

\subsection{Mapa reżimów \slid{conclusions}}

Na podstawie nocy A i P5 (noc B w toku):
\begin{enumerate}
  \item \textbf{Twirling przy równym budżecie} obniża błąd CDR przy 39 CZ o około 30\% (7/8 par); przy
  12 CZ nie ma efektu, bo błąd jest już na poziomie szumu. Wynik pilota (połowa błędu) jest więc
  prawdopodobnie napędzany głębszymi szczeblami; rozstrzygnie noc B.
  \item \textbf{ZNE} jest prawdziwe przy 12 CZ (lepsze od garbage), a przy 39 CZ jego zysk jest taki
  sam jak z garbage, czyli jest artefaktem zaniku sygnału. $W_{\text{neg}}$ i SNR tego nie wykryły,
  tylko kontrola garbage. ZNE jest za to około 42 razy tańsze od CDR.
  \item \textbf{CDR z twirlingiem} ma najmniejszy błąd na obu szczeblach; spread pomaga przy 12 CZ, przy
  39 CZ jeszcze nie.
  \item \textbf{Surowy błąd napędza czas obwodu} ($\times1{,}43$ na $\mu$s), nie sama liczba CZ; po CDR
  zostaje część od CZ ($\times1{,}34$ na 10 CZ).
\end{enumerate}
Dwie osie z teorii nadal obowiązują: wzdłuż budżetu próg $S^*=\kappa V/(B^2-B_{\text{ZNE}}^2)$,
wzdłuż głębokości zanik sygnału zamykający okno ZNE. O mechanizmie twirlingu: działa empirycznie, a
mechanizm wykracza poza prosty błąd $ZZ$ na CZ.

\subsection{Ograniczenia \slid{limitations}}

\textbf{Jedno urządzenie.} Granice zależą od $\gamma$ i od tego, ile CZ dokłada routing. Na kracie
(Garnet, Emerald) ta sama funkcja kosztuje inną liczbę CZ, a szum ma inną strukturę (przesłuchy z
wieloma sąsiadami). Liczby się przesuną, metoda wyznaczania granic pozostaje.
\textbf{Jedno zadanie.} Średnie $Z$ ma specyficzną podłogę ($T_1$) i ideały Clifforda w zbiorze
$\{-1,-0{,}8,\dots,1\}$, często zero. Inna obserwabla zmieni i $f$, i rozkład ideałów.
\textbf{CDR przy głębokich obwodach.} Degeneracja fitu to ograniczenie metody dla tej obserwabli,
nie naszej realizacji; motywuje spread.

\subsection{Co dalej \slid{next}}

Ten sam protokół na kratach IQM (Garnet 20, Emerald 54 kubity). Emerald jest za duży na symulację
($2^{54}$ amplitud), ale CDR nadal działa, bo ideały Clifforda liczy symulator stabilizatorowy.
Inne obwody i obserwable, w tym jednokubitowe, dla których ideał rzadziej jest zerem.

% =====================================================================
\clearpage
\appendix
\section*{Ściąga wzorów}
\addcontentsline{toc}{section}{Ściąga wzorów}
\begin{center}
\renewcommand{\arraystretch}{1.45}
\begin{tabular}{@{}p{3.6cm}p{11.4cm}@{}}
\toprule
model szumu & $E(\lambda)=f+(\Oid-f)e^{-\gamma\lambda}$, $B\approx-\gamma(\Oid-f)$ \\
Richardson & $\hat E=\sum c_kE_k$, $c_k=\prod_{j\ne k}\frac{\lambda_j}{\lambda_j-\lambda_k}$; $\{1,3,5\}$: $(\frac{15}8,-\frac54,\frac38)$ \\
obciążenie ZNE & $|B_{\text{ZNE}}|\lesssim2{,}5\gamma^3|\Oid-f|$ \\
$\kappa$ & równy $|\Lambda|\sum c_k^2-1$ (4 / 14,66 / 44,56); Neyman $(\sum|c_k|\sigma_k/\sigma_1)^2-1$ (3 / 11,25 / 35) \\
próg & $S^*=\kappa V/(B^2-B_{\text{ZNE}}^2)$ \\
zanik & $E(\lambda>1)=f\Rightarrow\hat E=c_1E_1+(1-c_1)f$ \\
$W_{\text{neg}}$, SNR & $\sum_s\max(0,-\sum_kc_kP_s(\lambda_k))\le0{,}02$; $|E(\lambda_{\max})|/\sqrt{V/N}\ge2$ \\
CDR & $\hat O=aE+b$, $a=\Cov/\Var$; EIV $a=\Cov/(\Var x-\bar\sigma_x^2)$, limit $4\times$ \\
twirling & $B\,\CZ\,A=\CZ$; $\Lambda_{\text{tw}}=\sum_\alpha\chi_{\alpha\alpha}P_\alpha\cdot P_\alpha$; koherentny $N^2\theta^2/4$ vs $N\theta^2/4$ \\
budżet & $(n_{\text{train}}+1)MS=1\,032\,192$ \\
ansatz & $[R_X(wx)\,R_YR_Z\,\text{splątanie}]^{3}\,R_YR_Z$, 55 wag, $f(x)=\sum_{\omega}c_\omega e^{i\omega x}$ \\
parameter shift & $\partial_\theta\ev O=\frac12[\ev O_{\theta+\pi/2}-\ev O_{\theta-\pi/2}]$ \\
McNemar & $p=2\sum_{j\le\min(b,c)}\binom{b+c}j2^{-(b+c)}$ \\
Holm & odrzucaj, dopóki $p_{(i)}\le\alpha/(m-i+1)$ \\
\bottomrule
\end{tabular}
\end{center}

\section*{Literatura}
\addcontentsline{toc}{section}{Literatura}
\small
\begin{itemize}[leftmargin=*]
  \item K.~Temme, S.~Bravyi, J.~Gambetta, PRL 119, 180509 (2017) --- ZNE.
  \item T.~Giurgica-Tiron i in., IEEE QCE (2020) --- cyfrowy folding.
  \item P.~Czarnik, A.~Arrasmith, P.~J.~Coles, L.~Cincio, Quantum 5, 592 (2021) --- CDR.
  \item J.~J.~Wallman, J.~Emerson, PRA 94, 052325 (2016) --- randomized compiling.
  \item R.~Takagi i in., npj Quantum Inf.\ 8, 114 (2022); Y.~Quek i in., Nat.\ Phys.\ (2024) --- granice kosztu QEM.
  \item A.~Kandala i in., Nature 549, 242 (2017) --- hardware-efficient ansatz.
  \item A.~Pérez-Salinas i in., Quantum 4, 226 (2020) --- data re-uploading.
  \item M.~Schuld, R.~Sweke, J.~J.~Meyer, PRA 103, 032430 (2021) --- modele kwantowe jako szeregi Fouriera.
  \item J.~Bowles, S.~Ahmed, M.~Schuld, arXiv:2403.07059 (2024) --- benchmarki QML.
  \item D.~Gottesman (1998); S.~Aaronson, D.~Gottesman, PRA 70, 052328 (2004) --- symulacja Cliffordów.
  \item Köster, Mauerer, arXiv:2607.09360 (2026) --- degeneracja ZNE.
  \item Scavino Alfaro, arXiv:2605.08251 (2026) --- $\kappa$ i alokacja Neymana.
  \item F.~Wilcoxon (1945); S.~Holm (1979); Q.~McNemar (1947).
\end{itemize}

\end{document}
